\section{Dynamic Adjustment and the Price of Commitment}
\label{sec:r11_integrated}
Why does a technology for changing preferences alter the initial financial position, and what does it cost to sustain the activity in which that technology is used? These questions require a common continuation economy. A preference change can improve the value of remaining in operation, accelerate access to discharge, or change the exposure to a future surrender option. A financial guarantee and a compulsory operating term can implement the same continuation allocation at different resource costs. We establish these distinctions before developing the reusable numerical operators that certify them.

\subsection{The Continuation Economy}
The state is $x=(u,X)$, with preference index $u$ and managed wealth $X$. An operating action is $a=(c,\theta,\pi)$: consumption, deliberate preference drift, and the risky share. The eight-date settlement economy, its cardinal preference primitive, and its boundary discharge rule are specified in Sections \ref{sec:preferences} and \ref{subsec:finite_protocol}. Operating benefits are denoted by $d$, the probability of the second transition regime by $\lambda$, and the elective surrender charge by $F$. Surrender pays $G(x)-F$ after a compulsory first interval; discharge at a contractual boundary pays $G(x)$. The guarantee and any signing payment use a separate utility numeraire and do not enter managed wealth.

At date zero, a positive share belongs to class $+$ and a nonpositive share to class $-$. This classification lasts for the entire first interval, not an isolated instant. The initial share is optimized continuously subject to long and short permissions $L$ and $S$, conditional on the inherited finite consumption--adjustment pairs and feasible frozen proposals. Later operating menus are unchanged. Section \ref{subsec:r9_permissions} gives the exact breakpoint representation. Thus the first-date permission experiment and the subsequent stopping problem are now evaluated as one target, rather than as separate favorable slices.

Write $W^r_\sigma$ for the initial class value in regime $r\in\{\mathrm{adj},0\}$, where $0$ prohibits deliberate adjustment but retains stochastic preference changes. Let $W^\uparrow_\sigma$ restrict deliberate adjustment to nonnegative values at every date and state. Define
\begin{equation}
 \Delta^r=W^r_+-W^r_-,\qquad
 \Omega_\sigma=W^{\mathrm{adj}}_\sigma-W^0_\sigma.
 \label{eq:r11_decisions}
\end{equation}
The relative option is $\Omega_+-\Omega_-=\Delta^{\mathrm{adj}}-\Delta^0$. This identity specifies the comparison; the transport inequalities and regional bounds below determine its sign.

\subsection{A Dynamic Transport Test}
For fixed economic primitives, write an operating backup as
$Q_n^a(v)=R_n^a+K_n^a v$, where $K_n^a$ is the nonnegative discounted live-state kernel. Boundary payments are included in $R_n^a$. Let $\mathcal R_n$ denote a set of reachable states. The sets include the support of every feasible first-date position under the largest permissions in the parameter region and are closed under every later feasible live transition. They are not the support of a selected optimal policy.

Consider a negative-drift action $a$ and a feasible zero-drift action $a^0$ with the same consumption and risky share. For endpoint regime $k\in\{0,1\}$, set
\begin{equation}
 D_{n,k}^a=K_{n,k}^a-K_{n,k}^{a^0}.
\end{equation}
Suppose $\ell_{n+1}\leq V^\uparrow_{n+1}\leq h_{n+1}$ throughout the parameter region, and put $v=(h+\ell)/2$, $e=(h-\ell)/2$. If the benefit interval is $[d_0,d_1]$, define
\begin{align}
 H_{n,k}^a={}&\max_{d\in\{d_0,d_1\}}
       \{R_{n,k,d}^a-R_{n,k,d}^{a^0}\}
       +D_{n,k}^a v+|D_{n,k}^a|e,\\
 \overline g_n^a={}&\max_{\lambda\in\{\lambda_0,\lambda_1\}}
       \{(1-\lambda)H_{n,0}^a+\lambda H_{n,1}^a\}
       +\varepsilon_n^a .
 \label{eq:r11_transport_test}
\end{align}
Absolute value is entrywise, and $\varepsilon_n^a$ is an evaluation-error allowance. The surrender charge enters the continuation bounds, although it is not paid by either current operating action.

\begin{proposition}[Reachable-state elimination of downward adjustment]
\label{prop:r11_transport}
Suppose the continuation enclosures and reachable sets just described are valid. Suppose every feasible negative-drift action has the indicated feasible zero-drift counterpart and $\overline g_n^a<0$ at every reachable live state. Then, simultaneously throughout the parameter region and for both initial position classes,
\begin{equation}
 W^{\mathrm{adj}}_\sigma=W^\uparrow_\sigma.
 \label{eq:r11_upward_equality}
\end{equation}
The conclusion holds at every date on its corresponding reachable set. It makes no assertion about alternative initial states outside those sets. A negative frozen proposal without a feasible counterpart must be checked separately; it cannot be omitted from the test.
\end{proposition}

This is a test on economic transition and reward primitives with certified continuation bounds, not a test on the sign of the reported initial control. The enclosures can be constructed without assuming an optimal selector: apply the upward-restricted Bellman recursion to the lower and upper continuation arrays, minimize or maximize over the endpoint laws, and use the adverse benefit and fee endpoints. Positivity of the kernels preserves the enclosure. The signed row difference in \eqref{eq:r11_transport_test} retains cancellation that a difference of two unrelated value-error bounds would discard.

The economic content can also be seen directly. When preference-boundary crossing is absent and the within-interval stopping fraction $\alpha$ is determined by wealth alone, changing drift from zero to $\theta$ shifts the arrival preference index by $\alpha h\theta$. Let $\mathcal V$ be interpolated continuation on a live branch and $G$ on a discharged branch. Apart from the quadratic adjustment cost, the change in the backup is
\begin{equation}
 \E\left[e^{-\rho\alpha h}
   \int_{u'_0}^{u'_0+\alpha h\theta}
          \partial_u\mathcal V(z,X')\,dz\right]
 -\frac{k\theta^2}{2}\E\left[\frac{1-e^{-\rho\alpha h}}{\rho}\right].
 \label{eq:r11_shadow_transport}
\end{equation}
Piecewise derivatives suffice for the settlement interpolation. A positive drift is profitable when its discounted continuation gain exceeds its resource cost. A negative drift may instead be valuable close to a preference-discharge boundary. Formula \eqref{eq:r11_transport_test} tests these effects without assuming that continuation is globally monotone or differentiable.

At the original central contract $(\lambda,d)=(0.125,0.425)$, both adjusted classes choose $(c,\theta)=(0.8,0.2)$. The local gain from this drift, holding the selected financial control fixed and using adjusted continuation, is $0.01143610$ in class $+$ and $0.01138571$ in class $-$. The common current adjustment cost is $0.00498752$. The larger part of the difference between the two total adjustment options is not that common cost: it is their different exposure to future adjustment opportunities. The relative local gain is $0.0000503943$, and the relative future-option exposure is $0.0003556734$, summing to $0.0004060676$ in the noncancellable economy. Supplement S.14 gives the exact decomposition and its selected-control reconstruction. With $F=0.80$, the positive unadjusted class gains an elective surrender option, reducing the relative future component while leaving the adjusted local comparison unchanged.

These are not the directions in the two-stage primitive example of Section \ref{sec:r7_primitive}. That example establishes a separate consumption-exposure mechanism and remains intact. In the dynamic economy, even the baseline cardinal flow can favor a higher preference index: at $c=0.8$, its derivative in $u$ is positive on $[1.2,2.8]$. The dynamic conclusion nevertheless requires continuation and discharge, not that flow derivative alone.

The distinction across initial states is substantial. At $(u,X)=(1.25,1.59375)$ and the central noncancellable contract, full adjustment exceeds upward-only adjustment in the nonpositive class by $1.22113089$. Downward-only adjustment exceeds no adjustment by $1.30273111$. Holding the full policy's initial financial controls fixed, negative drift increases the discounted next-date mass at the lower preference boundary from $0.17589502$ to $0.49750624$. This is arrival through the stipulated settlement lottery, not a claim about a diffusion hitting probability. The reduction of costly continuing exposure accounts for a one-step gain of approximately $0.43835$ before the adjustment and liquidation-benchmark terms. That initial state is outside the support tube certified below. There is therefore no global monotonicity claim hidden in \eqref{eq:r11_upward_equality}.

\subsection{A Joint Economic Domain}
\begin{theorem}[Adjustment, active surrender, and financial permissions]
\label{thm:r11_joint}
Consider the stored-array settlement target and the stored piecewise-affine first-date operator just specified, with $x_0=(2,1.25)$ and one compulsory interval. Let
\begin{equation}
\begin{split}
 \lambda&\in[0.12,0.13],\qquad d\in[0.424,0.4245],\qquad
 F\in[0.7999,0.8001],\\
 L&\in[0.7999,0.8001],\qquad S\in[0.4999,0.5001].
\end{split}
\label{eq:r11_box}
\end{equation}
Throughout this five-dimensional region, \eqref{eq:r11_upward_equality} holds at all reachable live states and dates. With outward-rounded displayed bounds, the initial comparisons satisfy
\begin{align}
 \Delta^{\mathrm{adj}}&\in[0.00018762,\;0.00020240],\label{eq:r11_adj_bound}\\
 \Delta^0&\in[-0.00017930,\;-0.00012218],\label{eq:r11_zero_bound}\\
 W^0_+(F,1)-W^0_+(0,8)&\geq0.000034708.
 \label{eq:r11_active_bound}
\end{align}
Consequently the relative adjustment option exceeds $0.0003098$. The positive-position comparison in the no-adjustment regime has a strictly valuable surrender margin: a policy in that class whose value is within $0.000034708$ of its supremum cannot avoid elective surrender almost surely.
\end{theorem}

The last statement concerns the constrained positive class. In this region the unconstrained no-adjustment agent prefers the nonpositive class; the theorem does not assert that this selected agent surrenders. An active counterfactual option changes the comparison of financial positions even when it is not exercised under the ultimately selected position. This distinction is essential to interpreting the contracting experiment.

The proof combines two certificates. The first establishes the transport inequalities on the complete feasible-support tube. Its reachable-node counts, including boundary nodes, are
$1,30,100,203,342,495,637,735,833$. Every negative operating drift is strictly dominated by its zero-drift counterpart; the least slack after its arithmetic allowance exceeds $0.000070902$. Frozen proposals have nonnegative drift at every reachable live node. The second certificate uses feasible-policy Bernstein coefficients and either the signed transition-law chord or the count-information upper value. It treats the benefit and fee corners jointly and uses the smallest permissions for feasible lower policies and the largest permissions for upper values. No interpolation of optimized values across the law parameter is asserted to be an upper bound without the transition correction.

The fee interval in \eqref{eq:r11_box} is not the earlier implementation interval $[0.85,0.90]$. The latter disables elective surrender throughout the original finite-menu law--benefit rectangle; it implements the noncancellable mandate. The new region instead establishes a strictly valuable surrender option in the indicated class and optimizes financial permissions jointly. Neither result absorbs the other. The unfavorable original corner at $F=0.80$, and the interventions $(L,S)=(0.82,0.50)$ and $(0.80,0.51)$, remain in the paper and are reoptimized in the new replication package. At the central contract, the first permission intervention makes both regimes prefer a positive position, whereas the second makes both prefer a nonpositive position. These are boundaries of the economic result, not errors to be removed from its record.

\subsection{Pricing the Choice of an Operating Institution}
\label{subsec:r11_institutions}
A financial guarantee should not be compared with a freely selectable compulsory term whose execution is left unpriced. We therefore specify both technologies. Guarantee capacity $E$ costs the agent $C(E)$ in the separate numeraire, with $C$ continuous and nondecreasing. A compulsory term $m$ costs the principal $D(m)$. These costs pay for different enforcement services. Fees accrue to the external enforcement sector, not to managed wealth; they are inactive in the implementation comparison below. No term cost was included in the original zero-cost procurement table, which remains a useful limiting case.

Fix an adjustment regime and an operating policy selected from the noncancellable optimum, including a specified selection among ties. Write $W^r=\max_\sigma W^r_\sigma$ for its agent value and $A_r$ for delivered discounted service. Let $F_r^*(m)$ be the statewise enforcement threshold in Proposition \ref{prop:r9_enforcement}, computed on the common feasible-support tube. The principal can choose any of the eight terms and any sufficient capacity, with $F=E$. At the threshold, continued operation is the stated selection among indifferent best responses. Strictly larger fees give the same implementation without that stopping tie.

\begin{proposition}[The priced implementation frontier]
\label{prop:r11_institution}
Under these primitives, the minimum cost of procuring the selected noncancellable allocation is
\begin{equation}
 P_r=\min_{1\leq m\leq8}
 \left\{[G(x_0)-W^r+C(F_r^*(m))]_++D(m)\right\}.
 \label{eq:r11_procurement}
\end{equation}
Procurement is feasible for a principal with service value $b$ if and only if $bA_r\geq P_r$. When participation binds for all eight alternatives, the chosen term minimizes $C(F_r^*(m))+D(m)$. The capacity frontier for adjustment is weakly below that for no adjustment, but comparisons of principal surplus must still charge the regimes' possibly different delivered service.

For a short-term financial implementation and a costless-capacity hard term, put $z=G(x_0)-W^r$. The short-term implementation is preferred precisely when its saving in direct term cost is at least
\begin{equation}
 [z+C(E)]_+-[z]_+.
 \label{eq:r11_institution_comparison}
\end{equation}
This expression equals $C(E)$ when participation binds and lies between zero and $C(E)$ otherwise. In particular, a freely selectable, zero-cost full term weakly dominates the positive-cost financial implementation; it strictly dominates when the additional grant is positive.
\end{proposition}

Equation \eqref{eq:r11_procurement} is an endogenous choice within an explicit enforcement technology, not an assertion of unrestricted optimal mechanism design. It also makes the original objection testable. At zero term cost, the model selects the free full term. With $C(E)=0.02E^2$ and $D(m)=\kappa(m-1)/8$, the central economy selects $m=8$ when $\kappa=0$, $m=7$ when $\kappa=0.01$, and $m=1$ when $\kappa=0.02$, in both adjustment regimes. The short-term capacities are approximately $0.77675422$ with adjustment and $0.81353318$ without it. All eight alternatives, grants, service amounts, and break-even service values are reported in the replication record; fees $10^{-5}$ above each threshold independently reproduce both implemented class values.

Thus preference adjustment has two distinct economic consequences. Its continuation incentives change the ranking of financial positions on \eqref{eq:r11_box}. Its improvement of operating value lowers the capacity required to implement a specified continuation allocation. The latter result becomes an institutional comparison only after both capacity and term execution are priced. An active surrender contract is not relabeled as an optimal implementation contract merely because it belongs to the same family.

\subsection{What the Computational Certificate Contributes}
The economic conclusion is a joint restriction on direction, financial choice, stopping, and the price of enforcing operation. A table of point solutions cannot establish the first three on a parameter continuum. Reward-feature transfer alone also cannot do so when $\lambda$ changes the transition law at every date. The signed correction of Theorem \ref{thm:corrected_chord} and the feasible-policy representation supply the missing law-parameter bounds; the transport test then uses statewise continuation enclosures to explain the dynamic direction rather than merely label a relative option.

The new chord and count runs use identical action targets, reward corners, and policy banks. Both certify the bounds above. Their work account charges the common kernel construction, endpoint solves, and validation to each method, while recording incremental certificate work separately. Earlier same-target precision frontiers, neural-free comparisons, structural quadratic controls, and unfavorable experiments are retained. This comparison supports a claim about the cost of certifying an economic decision, not the necessity of neural training or a universal computational advantage.

The numerical allowance is $10^{-7}$ per compared value; the difference bounds already include twice that amount. It encloses evaluation of the specified stored arrays, not the transcendental construction of another target. To transfer all three displayed value inequalities to another model, an additional common error for each compared initial value below $1.7354\times10^{-5}$ would suffice. This is a required budget, not an estimate of that model's error. Transferring the directional elimination additionally requires valid statewise reward and kernel bounds on its reachable sets. The diffusion formulation and its approximation analysis retain their separate role.
